\section{Introduction}
\label{sec:introduction}

\citet{lasser2026horizon} establishes the anti-monopolar
property within Goal-Frontier Maximization (GFM): under discounting,
full capability domination by a single
substrate is anti-maximizing when the discount factor $\gamma$ exceeds
a threshold $\gamma^* = 1 - \rext / \Delta_0$ (Proposition~6,
``Domination is Anti-Maximizing'').  The monopolar endpoint is the
structural analogue of the instrumental convergence
concern~\citep{omohundro2008basic, bostrom2014superintelligence}:
a sufficiently capable optimizer acquires resources and eliminates
competitors, but here the convergence is measured against a
structural objective ($\volL$) rather than inferred from
instrumental reasoning.  But
this result characterizes the \emph{endpoint}, not the \emph{path}.
A coalition that has reached effective substrate diversity $\meff = 1$
has already paid the anti-maximizing cost.  The question is whether the
path to that state is self-correcting (the coalition detects its error
through observation-channel feedback and reverses course) or
absorbing: each step makes the next step's error less detectable,
until the monopolar state becomes a fixed point from which no locally
rational trajectory escapes.

Three compound feedback loops share this structure.  They were
identified in the gap analysis across the GFM sequence and share a
common dynamical form---a locally rational decision narrows the
observation basis, making the \emph{next} locally rational decision's
harm less detectable---that is closely related to the ``what failure
looks like'' scenario of~\citet{christiano2019what}, where
AI systems gradually become less aligned through feedback-loop
degradation rather than through a single catastrophic misalignment.

\paragraph{Loop 1: Structural avoidance $\to$ leverage blindness $\to$ measurement neglect.}
The actor avoids structural discovery because learning a new
subsumption edge can contract $\volL$ in the short
term~\citep{lasser2026poset}.  This inflates the actor's leverage
estimates for capabilities whose subsumption structure is unknown,
reinforcing the avoidance.  The structural discovery
value~(Proposition~2 of~\citet{lasser2026horizon}) provides the
restoring force; the question is whether it acts quickly enough.

\paragraph{Loop 2: Partial subsumption $\to$ world-model homogenization $\to$ detection failure.}
Subsumption removes agents whose observation channels detected
specific kinds of miscalibration.  The actor's world model degrades on
those dimensions, making the next subsumption step appear more rational.
This is the primary case studied in detail in
Sections~\ref{sec:coupled_system}--\ref{sec:phase_boundary}.

\paragraph{Loop 3: Over-calibrated restriction $\to$ discovery avoidance $\to$ structural ignorance.}
Restriction calibrated above the value-positive threshold makes
structural discovery locally suboptimal, recreating the avoidance
pathology.  Risk estimates on restricted capabilities cannot be
updated (no exercise $\to$ no evidence), and under max-aggregation the
highest-risk estimate dominates~\citep{lasser2026horizon}.

\medskip

This paper formalizes the common dynamical structure underlying all
three loops and finds the phase boundary that separates
self-correcting trajectories from absorbing ones.

\subsection{Framing precondition: vol\_P as an operational target}
\label{sec:framing_precondition}

All results in this paper are stated in terms of $\volL$ dynamics and
$\volL$-error $\Lyap(\Wm_t)$.  We assume throughout that $\volL$ is an
operationally meaningful target: that is, improving $\volL$ corresponds
to an improvement in the experiential capability volume the framework
is designed to protect.  This is the precondition under which all the
bounds and convergence claims of this paper carry their safety
interpretation.  If $\volL$ diverges from the experiential target
(the \emph{B-to-C gap}~\citep{lasser2026gfm}), $\Lyap$ measures error
against the wrong quantity and every subsequent result inherits the
divergence.

This is not a minor limitation to be deferred.  The B-to-C gap is
the precondition for the framework itself, not a correction to be
layered on top.  A complete safety result would require establishing
that $\volL$ tracks the experiential target closely enough over the
operating regime that $\Lyap$-based bounds translate into
experiential-capability-volume bounds.  Establishing this is a
distinct research program (the realized-capability-volume
construction, in preparation); we import its assumption.  Readers who
doubt $\volL$ is operationally meaningful should read this paper's
results as conditional on that assumption holding, not as absolute
safety claims.

\subsection{Summary of results}

\begin{enumerate}
\item \textbf{The coupled dynamical system}
  (Section~\ref{sec:coupled_system}).
  We model the interaction between the capability poset $P_t$ and the
  actor's world model $\Wm_t$ as a coupled discrete-time system.
  Subsumption actions are locally rational under the (possibly
  miscalibrated) world model, and the world model's accuracy depends
  on the observation channels that survive in the current poset.

\item \textbf{The phase boundary theorem}
  (Section~\ref{sec:phase_boundary}).
  The central result: a characterized boundary in terms of the
  effective correction rate (exogenous channel signals $\rS$ plus
  endogenous calibration incentive $\alphamin \cdot c_{\mathrm{V}}
  \cdot \underline{\nu}$),
  the error-proportional degradation rate $\rW$, the subsumption
  frequency $\rsub$, the exogenous drift rate $\dmax$ from lost
  channels, and the minimum cross-substrate redundancy
  $\rhomincross$.  Below the boundary, the
  Lyapunov function $\Lyap(\Wm_t)$ converges to a neighborhood of
  zero whose radius shrinks as subsumption ceases, and once inside
  that neighborhood the anti-monopolar property applies.  Above it, the system enters an
  absorbing basin (under full blindness) or a metastable basin (under
  partial indirect evidence).

\item \textbf{Unified treatment of three feedback loops}
  (Section~\ref{sec:unified_loops}).
  All three loops map to the same $(P, \Wm)$ dynamical system with
  different effective parameters.  The Wamura problem is a degenerate
  phase boundary where $\rS = 0$ (Proposition~\ref{prop:wamura}).

\item \textbf{Minimum channel redundancy}
  (Section~\ref{sec:design_criterion}).
  The phase boundary implies a design criterion: the minimum
  cross-substrate redundancy per safety-relevant dimension needed to
  stay in the self-correcting basin (Theorem~\ref{thm:redundancy}).
  For two-substrate systems, each substrate's channels are singularly
  load-bearing (loss of one substrate's channel immediately violates
  the criterion whenever $\Kmincons > 1$ and leaves zero margin
  otherwise).
\end{enumerate}

\subsection{What this paper does not do}

The phase boundary is about poset dynamics, not experiential
interpretation: it does not resolve the B-to-C gap (the divergence
between the capability volume $\volL$ measures and the experiential
volume it is supposed to track).  The analysis assumes the actor
optimizes $\volL$; inner-outer alignment (mesa-optimization) is
orthogonal.  Observation channels are assumed honest throughout; the
enforcement mechanism for channel integrity is the cryptographic
commitment protocol of~\citet{lasser2026exo}.
That work ensures channels are honest; this paper ensures enough honest
channels survive.

\subsection{Dependencies on prior results}

Table~\ref{tab:dependencies} summarizes which results from the GFM sequence
this paper relies on and how they enter the analysis.

\begin{table}[h]
\centering
\small
\begin{tabular}{@{}lll@{}}
\toprule
\textbf{Source} & \textbf{Result used} & \textbf{Role here} \\
\midrule
Foundational & Self-trust $T_s$ (Appendix B) & Restoring force via prediction residuals \\
Poset        & Axiom verification (Prop.\ 1), leverage & Properties $\volL$ preserves or violates \\
Horizon      & Anti-monopolar property (Prop.\ 6) & Linearized growth model we extend \\
Horizon      & Structural discovery value (Prop.\ 2) & Self-correction force in loop 1 \\
Horizon      & Substrate physics (Defs.\ 10, 12--13) & Substrate partition for phase boundary \\
Causal       & Multi-channel architecture (Table~1) & Instances of observation channels \\
Causal       & Risk-trust $L^2$ convergence (Prop.\ 2) & Template for bounding correction rates \\
Verification & Channel isolation \& independence & Integrity conditions on channels \\
\bottomrule
\end{tabular}
\caption{Dependencies on prior GFM results.
  Foundational~=~\citet{lasser2026gfm},
  Poset~=~\citet{lasser2026poset},
  Horizon~=~\citet{lasser2026horizon},
  Causal~=~\citet{lasser2026scm},
  Verification~=~\citet{lasser2026exo}.}
\label{tab:dependencies}
\end{table}
